% FORMALISATION RATISS-FOCAL — Focalisation informationnelle
% Posé par Jonathan Evina, nuit du 21 sept. 2026. Compiler : pdflatex (ou Overleaf).
\documentclass[11pt,a4paper]{article}
\usepackage[french]{babel}
\usepackage{amsmath,amssymb}
\usepackage[margin=2.4cm]{geometry}
\title{RATISS-FOCAL : formalisation de la focalisation informationnelle}
\author{RATISS Labs --- Jonathan Evina}
\date{21 septembre 2026 (v2 : théorie + mesures)}
\begin{document}
\maketitle

\section{Question et lois accompagnatrices}
La cohérence émerge-t-elle de l'information ? Deux lois accompagnatrices
fournissent la mesure, sans machinerie neuronale.
\textbf{LCT} : $R = P_{\mathrm{sig}}$, $\Delta W = \eta\,\varphi\,P_{\mathrm{sig}}\,C$.
On certifie la forme, pas le courant.
\textbf{Stine-24} : soit $s_i(t)$ la part de persistance de la couche $i$,
$f_i(t) = 1 - |s_i(t) - s_i(0)|$ la fidélité de couche, alors
\begin{equation}
\Psi_{\mathrm{sync}}(t) = \left(f_1 f_2 f_3\right)^{1/3} \to 1 = \text{accordées}.
\end{equation}
Mesuré : $\Psi = 0.874$ (TEST-08, alignées).

\section{Échelles, couplage, intrication}
Trois couches (macro/micro/info), facteurs internes par couche plus facteur
externe $X$ :
\begin{equation}
X = \tfrac{1}{3}\sum_{i<j} |\mathrm{corr}(P_i(t), P_j(t))|.
\end{equation}
Mesuré : $X = 0.147$ (couplées, surtout macro-micro : $0.33$).
Non-séparabilité :
\begin{equation}
U = P_{\mathrm{joint}} - \sum_i P_i, \qquad |\Psi\rangle \neq |\mathrm{macro}\rangle \otimes |\mathrm{micro}\rangle \otimes |\mathrm{info}\rangle.
\end{equation}
Mesuré : $U_{\mathrm{fond}} = 0$ (contrôle), $U_{\mathrm{chargé}} = -1.36$ :
l'information lie les couches (TEST-10).
Tryperposition : collapse de 30\% d'une couche, sync préservée (ratio $> 0.7$).

\section{Influence et séparation}
Taux d'influence (terme d'interaction) :
\begin{equation}
I(X \to Y) = \frac{P_{Y+\mathrm{pont}} - P_Y - P_{\mathrm{pont}}}{A^X}.
\end{equation}
Mesuré : $I_{\mathrm{loin}} \approx -0.004 \simeq 0$, $I_{\mathrm{pont}} = +0.0025$ :
l'influence exige un contact. Borne $I_{\min} > 0$ : corrélation résiduelle
$0.26$ (bruits indépendants) contre $0.99$ (bruit commun) --- le couplage
réel porte la borne (TEST-12).
Séparateur global/individuel :
\begin{equation}
K(L) = 1 - \frac{P_L}{P_{\mathrm{joint}}}.
\end{equation}
Mesuré : $K = \{0.33, 0.80, 0.67\}$ (TEST-13).

\section{Secteurs et constantes}
Une seule loi ne suffit pas : résidus par secteur $/$ résidu global $= 0.001$
(secteurs $S_i$ requis, TEST-14).
Nombres sans dimension du labo :
\begin{equation}
\Pi_c = \frac{\Phi_c}{N\,\tau} = 0.0814, \quad
\Pi_P = \frac{P_{\mathrm{ref}}}{N_{\mathrm{pts}}} = 0.0230, \quad
\Pi_H = H = 0.997,
\end{equation}
stables entre tirages ($< 25\%$, TEST-15). Universalité : question ouverte.

\section{Concentration, porteurs, seuils}
\begin{equation}
\Phi(X) = \frac{N_{\mathrm{int}} \times \tau}{V_{\mathrm{info}}} \geq \Phi_c(X)
\;\Rightarrow\; \text{focalisation}.
\end{equation}
Porteurs : tranches disjointes $m_i \cap m_j = \varnothing$ vérifiées ;
franchissement $\Phi_c^{\mathrm{lab}} = 500$ aux étapes I$=3$, S$=4$ (TEST-03).
Seuils par structure : $\sigma_c(\text{boucle}) = 0.02 \neq
\sigma_c(\text{diffus}) = 0.20$ (TEST-16).

\section{Maintien, H2', correspondance}
F0 (focalisation dynamique) : entretenu $0.79$ contre abandonné $0.62$ sous
fort bruit, coût $1920$ bits (TEST-17).
H2' : $\kappa = f(D, I)$, grille $D \times I$ cartographiée, $R \in [0.97, 1.09]$ :
pas de bascule dans cette plage, ticket ouvert (TEST-18).
Correspondance : injection $\to 0$ = retour au témoin $P = 1.0$, monotone (TEST-19).

\section{Verdicts}
Anti-triche : dossier observé $/$ règles $= 2.03 > 2$ (émergence, de justesse).
Verdict : \textbf{A+B sœurs} (convergence émergence/implantation).
Univers-Unifié-v2 : 11/11, cohérent.
\end{document}
